\chapter[Sermon 97. The Conclusion of This Work, Wherein Is Established the Authority of the Same, and the Sum Collected Briefly.]{The Conclusion of This Work, Wherein Is Established the Authority of the Same, and the Sum Collected Briefly.}
\label{sermon:097}
\markright{Sermon 97. The Conclusion of This Work, Wherein Is Established ...}

And he said to me, “These sayings are faithful and true.” And the Lord God of the holy prophets sent his Angel to show to his servants the things which must be shortly fulfilled. Behold, I am coming shortly. Happy is he who keeps the prophecy of this book. I am John who saw these things, and I heard them. And when I had heard and seen, I fell down to worship before the feet of the Angel who showed me these things. And he said to me, “Do not do this.” For I am your fellow servant, and of your brethren the prophets, and of those who keep the sayings of this book. Worship God.

The sixth and last part of this work contains the conclusion, which affirms the things we have heard to be divine, certain, and undoubted: for he collects the chiefest things, and moves all men to faith, and study of godliness, that in steadfast hope we should look for the judge of all to come shortly, and to judge the living and the dead. And in goodly order this last book of the Canonical scripture finishes the godly narration and doctrine, with the judgment and end of all things. For the holy Scripture begins at the first origin of all things, and continues a narration until the end of all things, containing in itself the universality of things, and all such things as are requisite to be known about matters needful and profitable. And all those things our good Lord has given us to be known in the holy scripture, that is to say, in the Canonical books. For they are false harlots who say that all things which appertain to the true and full godliness and salvation of the faithful are not set forth in holy writings, and therefore have need of traditions. They indeed have need of those traditions, which will utter their crafty wares: we need none, who esteem all their wares not worth a farthing to be bought of any man. For Isaiah has sufficiently dissuaded us from their deceivable and crafty jugglings, in chapter 55. And this conclusion contains about 16 articles, which we shall discuss in order.

\index{Arethas of Caesarea}Immediately after the beginning is set a grave asseveration, that the things which he has said or written hitherto are true, sure, certain, and undoubted, he has in a manner the same sentence also in the nineteenth chapter of this book. And he calls faithful sayings, which are stable, ratified, and undoubted. And the sentence is referred in the things which he has spoken of the blessed life to the world to come, lest we should be left in any doubtfulness. Again it is referred to the whole narration of this book. And this sentence seems to be a clause of assertion, confirming the certainty of the matter proposed, as are those also in the prophets: “For the Lord has spoken,” and again, “Thus says the Lord of Hosts,” and that same most used in the gospel, “Verily, verily, I say unto you.” And in the apostolic Epistles, “God is my witness that I lie not.” And the goodness of God does succor our infirmity, whereby many times we, doubting of the verity of God’s words, do waver, and confirms our hope with these as it were anchors. Wherefore these must be diligently beaten in and urged in the ecclesiastical doctrine. Arethas expounding this place: as the wonted manner of this holy Evangelist is always, so is it here also. For like as in his gospel, in token of loyalty, he says, “And we know that his testimony is true,” so in this place also, setting to his seal, he says, “These sayings are faithful and true.” Hitherto he. Therefore it shall be an unworthy thing to doubt, be it ever so little, of the things that are written in this book, and in other books of the canonical Scripture.

\index{Antichrist}\index{John, Saint}Secondly, he repeats who is the author of this work, and that all these things are revealed to him. And truly, there is no other author but the Lord God himself, and that the God of the holy Prophets. This has great power, for he shows not only that he is one and the same God of both Testaments, who by his Spirit has inspired the prophets and Apostles; but also bids us secretly to esteem the truth and certainty of this book of the prophetic matters. For if he could in old time tell his people beforehand of things to come, and utter all things by the prophets, what marvel is it, if he now also performs the same by Saint John? And if all those things came to pass which the prophets did prophesy to come, neither did any word, no nor one jot fall to the ground which was not fulfilled, then there is no man also that will doubt of the truth of this book, if at least he consider that the same God which in times past was with the prophets, is now also with blessed John. The prophets said how the land of Canaan should be delivered into the possession of the children of Israel: it was delivered. The selfsame prophesied that the people of Israel should for their sins be cast out again of the same land into Babylon: they were cast out. After they prophesied again that they should be delivered, should repair the City, to which Christ would come, which should redeem mankind, and call into fellowship of life and bless all nations: they were delivered, they repaired their City: Christ came, and redeemed mankind, and the gospel was preached throughout the whole world. What thing then remains, but that the church should be troubled, Antichrist should come, and reign, and that the true Christians and he should wage battle together, and the Judge come at the last unto judgment, and reward every one according to his doings? And this place proves the divinity of Christ infallibly. For what can be spoken more plainly than was said? The Lord God of holy Prophets sent forth his Angel. So in the first chapter is said: The revelation of Jesus Christ, which God gave him. And a little after he says: I Jesus sent my Angel, that he might testify unto you, \&c. Herein therefore is shown the unity of the divine substance, and the destruction of persons.

And the manner of the revelation is shown, or repeated and collected rather: he sent his angel. Christ therefore, by his angel, shows all things to John. For no one has seen God at any time; neither shall the Lord come down again from heaven before the judgment. Wherefore this whole vision was exhibited and declared by the angel, which was the messenger of Christ the Lord. Wherefore all things are properly referred to Christ, who sent the angel. But to whom did he show or reveal these things? To his servant. For the scorners of God laugh at these things, and take them for fables. But God loves his worshippers, and warns them of all things in due season.

\index{Judgment, Last}Now he gathers together some of the things he has treated hitherto. These are chiefly contained in two points. For he shows what must be done, briefly. For this book contains the destinies of the church from the time of the Apostles to the end of the world. Therefore, he did not prophesy far off, but of things that began in the very time of Saint John. And if these things must be done, who shall resist? Not that I will establish the necessity of the Stoics, but that I acknowledge the mighty working of God, according to his providence and righteousness. After this, he adds another point: “Behold, I come quickly.” For this book comprehends many things which concern the judgment itself, and the last judgment, to which I will come so swiftly and unexpectedly that the wicked and ungodly shall expect nothing less. For the Lord says in the Gospel: “It shall be as in the days of Noah and Lot.” And in the hour that you think not, the Son of Man will come. Moreover, as brightness comes forth from the East and shines in the West, so shall be the coming of the Son of Man. And therefore, the Lord says now also at this present: “Behold, I come quickly.” For suddenly, while he seems in the meantime to do another thing, at unawares he brings in the Lord speaking, and that a wonderful matter, as this particle indicates. For Saint Paul has also written, while they shall say peace and security, sudden destruction shall come upon them.

4. But what benefit should God’s servants expect from this book? He encompasses much in a few words, and says: “Happy is he who keeps the words of the prophecy of this book.” Felicity and blessedness are the fruits to be gained from this book. In this present world, being united with Christ, we shall walk in the way of righteousness, and avoid the schemes of Antichrist; and shall not feel the torments that arise in the conscience, from the corruption of depraved religion. And when we depart from here, we shall go directly to those blessed places. This is the ultimate blessedness and felicity. And let us note that it is not enough either to have seen, or heard, or read this book; it must needs be kept. For we must beware that it does not enter one ear and exit the other, so that we forget the things that are told us, but rather that we shape our entire life according to the doctrine of this book. And he attributed to it the title of prophecy. All Scripture is called prophecy, meaning divine. But considering how this book, for the most part, reveals things to come for the church, it is rightly called a prophecy.

5. He repeats again and emphasizes, [?beats upon], both his name and also that he is a witness who saw and heard, who may surely be credited. And thus he will gain authority for this book. For it must needs be held in great estimation that which was conceived and written by the Apostle and Evangelist John. Many consider it a fault that John so diligently expresses his name. But it is marvelous that they do not observe the same elsewhere, and in the writings of others not without praise. Does not the same John repeat and impress the name of a disciple in the story of the Gospel? Who should reprehend this? I see not therefore what he has offended herein: but rather he foresaw in the spirit that many would speak against this book, not without great cause, and with much fruit. And also of extreme necessity he importuned his name. And the Apostle Paul also to the Galatians: “Behold, I Paul say unto you,” says he, “if you be circumcised, Christ shall profit you nothing.” The same also, to move affection, impresses his name to Philemon. Arethas therefore very aptly expounds this place. And this, says he, is a certain propriety of speech in this Apostolic soul. For even as he did in the Gospel also, where he says: “And he that saw hath borne witness, and his testimony is true,” the same doth he in this place also, testifying that he was both a hearer and beholder of these things which are prophesied. For hereby he wins credit to the things which had been seen. Thus much he. Others have thought that not without cause, John has in this book repeated his name oftener than in his story, for that men will more hardly believe a prophecy speaking of things yet to come, than a story which tells of matters past.

\index{Aquinas, Thomas}\index{Papacy}6. In the sixth place, he adds what happened to him again with the angel, revealing to him these great mysteries. A similar story would apply to all the world, as we have in the nineteenth chapter, where we also explained it; and anyone who wishes may see. And yet, the commentators ask how it happens that John again does what he did before and was prohibited by the angel? Thomas Aquinas suggests that John, being beside himself because of the splendor of the visions, did this as one who is astonished. The gloss says that before, the angel forbade him from worshipping with \textit{latria}, but here he prohibits him from worshipping with \textit{dulia}. But to me it appears—preferring always the better judgment of others—that in John is shown to all the godly how great is the frailty of humankind to fall, unless it is restrained and drawn back by the mighty hand of God. The angel had shown John expressly before that he should not do what he now did, and now repeats it again. For having, as it were, forgotten those things because of the splendor of the angel, he would surely have given him some worship. For so we permit to ourselves more than is proper, especially toward noble persons, whom for their excellent gifts of God we deem worthy, whom we may also without offense to God even worship. That opinion deceives most people in our time, those who against the comeliness of sincere religion worship and honor saints. But the angel of the Lord here neither invents nor brings forth any new doctrine, but that old in form, as they term it, so that we should understand that the will of God is always one and perpetual, which will does not want the most excellent creatures to be worshipped, but one God alone to be honored. He repeats therefore the same reasons, which he also objected before. Therefore be they always of force, with all, and at all times. John meanwhile seems to want to commend to us the splendor of this vision or revelation; and that the angel constantly admonished him of his duty, and us all by him, that the thing which is proper to God, we should not transfer to any creatures, and it deserves exceeding great praise here that John here conceals nothing, but by express words commits to writing his fall and rebuking of the angel most evidently. For by his fall he would admonish that the godly should not fall in like cases, but give all glory to God. Here seems also to be observed a marvelous affection in the manner of speaking. For the angel cries out to John, being ready to fall down now, yea prostrate already, and now about to worship: See thou do it not, that thou verily intend to do. Here is expressed the carefulness of mind and haste with which he goes about to prevent the enterprise of John. And thus diligent are the Holy Spirits in heaven in letting all things, that by any means do turn us from God, to the worshipping of creatures, much less would they themselves be worshipped, or have the things attributed to them, which the Papists at this day attribute by force of arms. The Lord of clemency and mercy convert them to a right mind, that they may attribute all glory to God.
